\section{Equal balances after unequal payment histories}
\label{sim:settlement-remedies}

A borrower can recover from a cash shortage before the reporting date. Its
closing statements then need not reveal that it failed to pay on time. To
isolate the timing, consider a synthetic company with cash 60, inventory 100,
debt carrying amount 80 and book equity 80. The contractual promise is 100,
payable on 30 June 2027. Assume a zero coupon, no fees or taxes, and one
remaining measurement interval that accretes the discount of 20 at maturity.
This deliberately simple interval makes the cash and contract comparisons
exact; a general effective-interest schedule needs its own measurement rules.

The company issues 40 new shares for one unit of cash each. In the first
history, the subscription settles on 29 June. Cash rises from 60 to 100,
the discount accretes on 30 June, and the borrower pays the complete promise
that day. In the second history, the shares settle on 2 July. The borrower
can pay only 60 at maturity. On 1 July an explicit observation records the
unpaid 40; the new cash arrives on 2 July and pays the remaining claim on
3 July. No overdraft, emergency loan or retroactive receipt funds the earlier
payment.

The following comparison follows the money and claim in the second history.
The liability column is positive for presentation, although journal credits
use the implementation's signed convention.
\begin{center}
\begin{tabular}{lrrr}
\toprule
After the dated event & Cash & Claim unpaid & Historical miss\\
\midrule
30 June: accretion &60&100&0\\
30 June: partial payment &0&40&0\\
1 July: missed-payment observation &0&40&40\\
2 July: share proceeds &40&40&40\\
3 July: final payment &0&0&40\\
\bottomrule
\end{tabular}
\end{center}
The zero in the first two historical-miss cells records that the observation
has not yet been made. It is not a finding that a payment was timely. The
observation changes contract memory without changing cash or posting another
expense. The final payment clears the current claim; it leaves the history
of the earlier shortfall intact.

Both histories end with the same financial statements under the supplied
treatment. Inventory of 100 is the only remaining asset, the debt is zero,
and equity is opening equity plus the subscription less the discount expense:
\begin{equation}
 A=100,\qquad L=0,\qquad E=80+40-20=100.
\end{equation}
Profit for the illustrated movements is $-20$. Net financing cash flow is
$40-100=-60$ in either history. The late path divides the repayment into
$60+40$, which changes its timing but not its total. Even agreement across
the closing balance sheet, income and cash-flow statements does not answer
whether the maturity was paid on time. A simulation query about missed
payments must retain the distinguishing history.

The difference also explains why balancing the journal is insufficient. Share
proceeds can debit cash and credit capital correctly while a faulty method
fails to increase the recorded shares. A missed-payment observation can omit
its history update without unbalancing any account. In the checked reference,
both omissions pass generic accounting execution and fail the independent
complete-state comparison. Explicit financing conditions subsequently reject
these proposed transitions before accepting their successors.

Those conditions have two roles. Snapshot conditions constrain current
quantities: claims and arrears are nonnegative, arrears cannot exceed the
remaining claim, and under this matured zero-coupon convention the carrying
liability equals that claim. Transition conditions constrain how an event
changes them. Payment of $p$ gives remaining claim $q'=q-p$, reduces cash by
$p$, and leaves the historical missed amount unchanged. A share subscription
must increase both the monetary capital account and the separate share
quantity under the supplied issue terms. An internally plausible rewrite of
historical miss from 40 to 50 can satisfy snapshot inequalities but violates
the transition condition. This is why both kinds of condition are needed.

Opening history creates another distinction. A supplied fiscal opening must
satisfy the declared period-start, temporary-balance and flow restrictions;
an intermediate state containing this year's payments is not such an opening.
Importing a consistent snapshot does not establish that it includes every
earlier missed payment. The selected opening method therefore compares it
with a separately supplied history premise and preserves unknown completeness.
Actual source evidence is still needed to establish that premise.

\section{A missed payment and the available remedies}
The amount unpaid does not by itself identify the creditor's remedy. Apple's
2024 annual report provides a concrete historical description of why the
distinction matters. Exhibit 4.1 expressly calls its description incomplete
and qualifies it by the 2013 and 2018 indentures; its issuer references concern
Apple Inc., rather than automatically every subsidiary
\citep[Exhibit 4.1, printed p.~3]{apple2024}. We use the summary to identify
conditions a simulation must represent. It does not establish an actual
issuer default, complete operative terms, or the applicability of those terms
to the synthetic company.

The described interest-payment event requires default continuing for 30 days
after the installment becomes due. Principal and premium use a different
condition: nonpayment when due and payable, including maturity, redemption
or declaration. Importing the interest grace period into principal repayment
would therefore apply the wrong described condition. The other-covenant route
has its own exclusions, a continuing breach and 90 days after qualifying
written notice. Holder notice uses at least 25\% of outstanding principal
under the 2013 indenture and at least 33\% under the 2018 indenture; the
trustee has a separately described notice route
\citep[Exhibit 4.1, printed p.~9]{apple2024}.

The thresholds concern the relevant outstanding series, not a fraction of
consolidated debt or a count of creditors. Series mapping, denominator,
delivery to the required recipients and validity of notice are substantive
inputs. No notice and no evidence about notice lead to different answers.
Likewise, day-level data at exactly the 30- or 90-day boundary may leave
event order unresolved. The conditional reference retains that uncertainty;
it does not assert that the complete contract has no answer at the boundary.

The summary next distinguishes declaration from automatic acceleration.
Certain continuing nonbankruptcy events permit a declaration by the specified
trustee or holder route. Certain bankruptcy events make the obligations due
without that declaration. A negative equity balance establishes neither the
specified bankruptcy event nor the required category. Moreover, an event for
one series does not necessarily constitute an event for another. Neither
automatic contagion nor universal independence follows
\citep[Exhibit 4.1, printed pp.~9--10]{apple2024}.

Cure and rescission must also remain separate. The described rescission of a
declaration requires holder action and specified deposits and cure or waiver
conditions, with an exception concerning amounts due solely because of
acceleration. Payment can remove current arrears without by itself proving
that an earlier acceleration has been rescinded. In the selected synthetic
continuation, payment changes the economic claim, reassessment changes the
current clause finding, and a separate supported rescission assessment changes
the recorded acceleration state. The historical miss remains 40 throughout.

Consent cannot be supplied by a numerical convenience either. The summary
requires each affected holder's consent for specified protected changes,
including changes to stated maturity and reductions of principal or interest.
Its general majority-waiver provision excludes payment defaults and certain
individually protected provisions. A waiver concerns the specified past
default; it does not license a future one
\citep[Exhibit 4.1, printed pp.~11--12]{apple2024}. A current contractual default
may cease under a valid waiver while the economic record still shows that
payment arrived late. Keeping those two statements is consistent.

The surrounding qualifications matter economically. Individual payment
enforcement differs from collective proceedings; legal discharge differs
from covenant defeasance. The description warns that assets sufficient to
fund scheduled maturity may be insufficient after acceleration
\citep[Exhibit 4.1, printed pp.~10--13]{apple2024}. A simulator must therefore
retain the relevant rights, legal owner, timing and usable funding, even
when a consolidated closing statement appears fully funded.

The typed conditional implementation translates the retained examples and
records selected findings through the common monetary and memory engine.
An unresolved finding produces no successor state. Recording a known finding
changes no cash, carrying amount, shares or historical missed amount. This
establishes a limited representation and execution result. Multiple concurrent
clauses, the full instruments, actual notices and amendments, and the economic
consequences of acceleration still require their own admitted methods.
